% ---------------------------------------------------------------------------
\begin{frame}{Third Normal Form (3NF)}
  \begin{block}{Definition}
    A relation is in \textbf{3NF} if it is in 2NF and no non-prime attribute is
    \textbf{transitively dependent} on the primary key.  Every non-prime
    attribute must depend \emph{directly} on the PK, not through another
    non-prime attribute.
  \end{block}

  \begin{itemize}
    \item \textbf{Transitive dependency:}
          $PK \to A \to B$, where $A$ is non-prime.  $B$ does not depend
          directly on the PK; it reaches the PK via $A$.
    \item In our \texttt{Albums} table:\\
          $\texttt{AlbumID} \to \texttt{MainArtist} \to \texttt{Debut}$\\
          The artist's debut year depends on the artist, not on the album.
    \item \textbf{Fix:} Extract a separate \texttt{Artists} table:
          \begin{itemize}
            \item \texttt{Albums(\underline{AlbumID}, Title, ArtistID, Genre, Release)}
            \item \texttt{Artists(\underline{ArtistID}, Name, Debut)}
            \item \texttt{Songs(\underline{AlbumID, Song}, FeaturedArtist)}
          \end{itemize}
  \end{itemize}

  \begin{block}{Why 3NF Matters in Practice}
    Most production databases are designed to 3NF.  It eliminates the most
    common anomalies with reasonable decomposition effort.  The
    \textbf{Synthesis Algorithm} always produces a 3NF schema that is both
    lossless and dependency-preserving.
  \end{block}
\end{frame}

% ---------------------------------------------------------------------------
\begin{frame}{Boyce-Codd Normal Form (BCNF)}
  \begin{block}{Definition}
    A relation is in \textbf{BCNF} (sometimes called \emph{3.5NF}) if it is in
    3NF and for \emph{every} non-trivial FD $A \to B$, $A$ is a
    \textbf{superkey}.
  \end{block}

  \begin{itemize}
    \item BCNF is \textbf{strictly stronger} than 3NF: it catches edge cases
          where 3NF still permits anomalies.
    \item In 3NF, the right-hand side of an FD is allowed to be a prime
          attribute; BCNF \emph{disallows} this exception.
    \item \textbf{Example violation:}
          \texttt{Releases(\underline{ID, Edition}, LabelID, Label)}\\
          FDs: $\texttt{LabelID} \to \texttt{Label}$ and
          $\texttt{Label} \to \texttt{LabelID}$.\\
          Neither \texttt{LabelID} nor \texttt{Label} is a superkey, yet they
          determine each other --- a BCNF violation.
    \item \textbf{Fix:} Decompose into:\\
          \texttt{Releases(\underline{ID, Edition}, LabelID)} and
          \texttt{Labels(\underline{LabelID}, Label)}
  \end{itemize}

  \begin{alertblock}{Trade-off}
    BCNF \emph{may not preserve all FDs} after decomposition.  If
    dependency-preservation is critical, 3NF (via the Synthesis Algorithm) may
    be preferred.
  \end{alertblock}
\end{frame}

% ---------------------------------------------------------------------------
\begin{frame}{BCNF -- Example}
  \textbf{Before (not in BCNF):}
  \vspace{4pt}

  \small
  \begin{tabular}{llll}
    \hline
    \textbf{ID} & \textbf{Edition} & \textbf{LabelID} & \textbf{Label} \\
    \hline
    A100 & 3am Edition        & REP & Republic Records  \\
    A101 & Late Night Edition & UMG & Universal Music   \\
    \hline
  \end{tabular}
  \normalsize

  \vspace{8pt}
  FDs: $\{\texttt{ID, Edition}\} \to \texttt{LabelID} \to \texttt{Label}$
  and $\texttt{Label} \to \texttt{LabelID}$.\\
  \texttt{LabelID} is \textbf{not} a superkey, yet $\texttt{LabelID} \to
  \texttt{Label}$ --- \textbf{BCNF violation}!

  \vspace{8pt}
  \textbf{After (BCNF):}
  \vspace{4pt}

  \small
  \begin{columns}
    \begin{column}{0.45\textwidth}
      \begin{tabular}{lll}
        \hline
        \multicolumn{3}{c}{\textbf{Releases}} \\
        \hline
        \textbf{ID} & \textbf{Edition} & \textbf{LabelID} \\
        \hline
        A100 & 3am Edition        & REP \\
        A101 & Late Night Edition & UMG \\
        \hline
      \end{tabular}
    \end{column}
    \begin{column}{0.35\textwidth}
      \begin{tabular}{ll}
        \hline
        \multicolumn{2}{c}{\textbf{Labels}} \\
        \hline
        \textbf{LabelID} & \textbf{Label} \\
        \hline
        REP & Republic Records \\
        UMG & Universal Music  \\
        \hline
      \end{tabular}
    \end{column}
  \end{columns}
  \normalsize

  \vspace{6pt}
  Now every determinant (\texttt{ID}+\texttt{Edition} and \texttt{LabelID}) is
  a superkey. BCNF satisfied.
\end{frame}
